\SourceLeaf{417}


\Needspace{9\baselineskip}\begin{center}{\Large SEVENTH SECTION}\par\medskip ON\par\medskip {\large THE EQUATIONS DEFINING DIVISIONS OF THE CIRCLE.}\end{center}
\BookRule{0.5\textwidth}

\Needspace{6\baselineskip}\Article{335}
Among the most brilliant advances added to mathematics by the work of recent mathematicians, the theory of functions depending on the circle undoubtedly holds a particularly distinguished place. The greatest mathematicians of recent times have devoted their industry and insight so persistently to this remarkable kind of quantity, to which we are so often led in investigations of the most varied kinds, and whose help no part of mathematics can do without, and have built from it so vast a subject, that one could hardly have expected any part of this theory, much less an elementary part lying, as it were, at its very threshold, still to admit substantial advances. I am speaking of the theory of trigonometric functions corresponding to arcs commensurable with the circumference, or the theory of regular polygons; this section will show how little of it has so far been developed. Readers might wonder why such an investigation is undertaken in this work in particular, devoted chiefly to a subject which at first sight seems entirely different; but the treatment itself will make sufficiently clear how closely this subject is connected with higher arithmetic.

The principles of the theory we are about to present also extend much further than we carry them here. They can be applied not only to circular functions, but with equal\SourceLeaf{418}
success to many other transcendental functions, for example those depending on the integral \(\int\frac{dx}{\surd{(1-x^4)}}\), and also to various kinds of congruences. But since we are preparing an extensive separate work on those transcendental functions, and congruences will be treated in detail in the continuation of the arithmetical investigations, we have thought it appropriate to consider only circular functions here. Indeed, even these, which we could treat in the fullest generality, we will reduce to the simplest case by the means to be explained in the following articles, both for brevity and to make the entirely new principles of this theory easier to understand.

\Needspace{9\baselineskip}\begin{center}\textit{The investigation is reduced to the simplest case, in which the number of parts into which the circle is to be divided} \textit{is a prime.}\end{center}


\Needspace{6\baselineskip}\Article{336}
Denote the circumference of the circle, or four right angles, by \(P\), and suppose \(m,\allowbreak n\) are integers, and \(n\) the product of pairwise coprime factors \(a,\allowbreak b,\allowbreak c\), etc. By Article 310 the angle \(A=\frac{mP}{n}\) can be put in the form \(A=(\frac\alpha a+\frac\beta b+\frac\gamma c+\mathrm{etc.})P\), and its trigonometric functions can be derived by known methods from the functions belonging to the parts \(\frac{\alpha P}a,\allowbreak \frac{\beta P}b\), etc. Thus, since we can take primes or powers of primes for \(a,\allowbreak b,\allowbreak c\), etc., it clearly suffices to consider division of the circle into a number of parts equal to a prime or a power of a prime, and the polygon with \(n\) sides will follow immediately from the polygons with \(a,\allowbreak b,\allowbreak c\), etc., sides. Here, however, we will restrict the investigation to the case in which the circle is to be divided into a number of parts which is an (odd) prime, chiefly for the following reason. It is known that the circular functions corresponding to the angle \(\frac{mP}{pp}\) can be derived from the functions belonging to \(\frac{mP}p\) by solving an equation of degree \(p\), and likewise from these, by an equation of the same degree, the functions belonging to \(\frac{mP}{p^3}\), etc.; thus, if the polygon with \(p\) sides is already given, determining the polygon with \(p^\lambda\) sides necessarily requires solving \(\lambda-1\) equations of degree \(p\). Although the following theory can also be extended to this case, this route will still lead to just as many equations of degree \(p\), which, if \(p\) is a prime, cannot in any way be reduced to lower degrees. For example, it will be shown below that the polygon with \(17\) sides can be constructed geometrically; but in determining the polygon with \(289\) sides there is no way to avoid an equation of degree \(17\).
\SourceLeaf{419}

\Needspace{9\baselineskip}\begin{center}\textit{Equations for the trigonometric functions of arcs which are one or more parts of the whole circumference;} \textit{reduction of trigonometric functions to the roots of the equation \(x^n-1=0\).}\end{center}


\Needspace{6\baselineskip}\Article{337}
It is well known that the trigonometric functions of all the angles \(\frac{kP}n\), where \(k\) ranges over all the numbers \(0,\allowbreak 1,\allowbreak 2\ldots n-1\), can be expressed by roots of equations of degree \(n\): namely, the \textit{sines} by the roots of (I) \[\begin{aligned}
x^n-\frac14nx^{n-2}+\frac1{16}\frac{n.n-3}{1.2}x^{n-4}
&-\frac1{64}\frac{n.n-4.n-5}{1.2.3}x^{n-6}\\
&+\mathrm{etc.}\pm\frac1{2^{n-1}}nx=0
\end{aligned}\] the \textit{cosines} by the roots of (II) \[\begin{aligned}
x^n-\frac14nx^{n-2}+\frac1{16}\frac{n.n-3}{1.2}x^{n-4}
&-\frac1{64}\frac{n.n-4.n-5}{1.2.3}x^{n-6}\\
&+\mathrm{etc.}\pm\frac1{2^{n-1}}nx-\frac1{2^{n-1}}=0
\end{aligned}\] and finally the \textit{tangents} by the roots of (III) \[\begin{aligned}
x^n-\frac{n.n-1}{1.2}x^{n-2}&+\frac{n.n-1.n-2.n-3}{1.2.3.4}x^{n-4}\\
&-\mathrm{etc.}\pm nx=0
\end{aligned}\] These equations (which hold generally for every odd value of \(n\), and II also for even values) can easily be reduced, on setting \(n=2m+1\), to degree \(m\): I and III by dividing the left side by \(x\) and substituting \(y\) for \(xx\). Equation II clearly has the root \(x=1\ (=\cos0)\), and the remaining roots are always equal in pairs (\(\cos\frac Pn=\cos\frac{(n-1)P}n,\allowbreak \ \cos\frac{2P}n=\cos\frac{(n-2)P}n\), etc.); thus its left side is divisible by \(x-1\), and the quotient will be a square. Taking its square root reduces equation II to \[\begin{aligned}
x^m+\frac12x^{m-1}-\frac14(m-1)x^{m-2}&-\frac18(m-2)x^{m-3}\\
&+\frac1{16}\frac{m-2.m-3}{1.2}x^{m-4}\\
&+\frac1{32}\frac{m-3.m-4}{1.2}x^{m-5}-\mathrm{etc.}=0
\end{aligned}\] whose roots will be the cosines of the angles \(\frac Pn,\allowbreak \frac{2P}n,\allowbreak \frac{3P}n\ldots\frac{mP}n\). Further reductions of these equations, at least for the case in which \(n\) is a prime, have not so far been available.

None of these equations, however, is as easy to handle or as suitable for our purpose as \(x^n-1=0\), whose roots are known to be very closely connected with theirs. Namely, writing \(i\) for the imaginary quantity \(\surd{-1}\) for brevity, the roots of the equation \(x^n-1=0\) are given by \[\cos\frac{kP}n+i\sin\frac{kP}n=r,\]
\SourceLeaf{420}
where \(k\) is to take all the values \(0,\allowbreak 1,\allowbreak 2\ldots n-1\). Thus, since \(\frac1r=\cos\frac{kP}n-i\sin\frac{kP}n\), the roots of equation I are given by \(\frac1{2i}(r-\frac1r)\), or by \(i\frac{1-rr}{2r}\); the roots of equation II by \(\frac12(r+\frac1r)=\frac{1+rr}{2r}\); and finally the roots of equation III by \(\frac{i(1-rr)}{1+rr}\). For this reason we will base the investigation on considering the equation \(x^n-1=0\), assuming \(n\) to be an odd prime. To avoid having to interrupt the order of the investigations, we give the following lemma first.

\Needspace{6\baselineskip}\Article{338}
\textsc{Problem.} \textit{Given the equation} \[(W)\ldots z^m+Az^{m-1}+\mathrm{etc.}=0,\] \textit{find an equation \((W')\) whose roots are the \(\lambda\)-th powers of the roots of the equation \((W)\),} \textit{where \(\lambda\) denotes a given positive integer exponent.} \textit{Solution.} Denote the roots of the equation \(W\) by \(a,\allowbreak b,\allowbreak c\), etc.; the roots of the equation \(W'\) must be \(a^\lambda,\allowbreak b^\lambda,\allowbreak c^\lambda\), etc. By Newton's well-known theorem, the coefficients of the equation \(W\) allow us to find the sums of any powers of the roots \(a,\allowbreak b,\allowbreak c\), etc. Thus find the sums \[\begin{gathered}
a^\lambda+b^\lambda+c^\lambda+\mathrm{etc.},\quad
a^{2\lambda}+b^{2\lambda}+c^{2\lambda}\ \mathrm{etc.}\ \mathrm{etc.}\\
\text{up to}\quad a^{m\lambda}+b^{m\lambda}+c^{m\lambda}+\mathrm{etc.}
\end{gathered}\] from which, proceeding in reverse by the same theorem, we can deduce the coefficients of the equation \(W'\). \textit{Q. E. F.} — At the same time this shows that, if all the coefficients in \(W\) are rational, all those in \(W'\) will also be rational. It can indeed be proved by another route that, if the former are all integers, the latter will also all be integers; but we will not dwell here on this theorem, which is not so necessary for our purpose.

\Needspace{6\baselineskip}\Article{339}
The equation \(x^n-1=0\) (under the assumption, always to be understood from now on, that \(n\) is an odd prime) has a single real root, \(x=1\); the remaining \(n-1\) roots satisfy the equation \[x^{n-1}+x^{n-2}+\mathrm{etc.}+x+1=0\] and are all imaginary. We will denote their collection by \(\Omega\), and the function \[x^{n-1}+x^{n-2}+\mathrm{etc.}+x+1\quad\text{by}\quad X\]
\SourceLeaf{421}
respectively. Thus, if \(r\) is any root in \(\Omega\), we have \(1=r^n=r^{2n}\), etc., and generally \(r^{en}=1\) for every integer value of \(e\), positive or negative; hence, if \(\lambda,\allowbreak \mu\) are integers congruent modulo \(n\), we clearly have \(r^\lambda=r^\mu\). But if \(\lambda,\allowbreak \mu\) are incongruent modulo \(n\), \(r^\lambda\) and \(r^\mu\) will be unequal; for in this case an integer \(\nu\) can be chosen so that \((\lambda-\mu)\nu\equiv1\ (\mathrm{mod.}\ n)\), whence \(r^{(\lambda-\mu)\nu}=r\), and therefore \(r^{\lambda-\mu}\) is certainly not \(=1\). It is also clear that every power of \(r\) is a root of the equation \(x^n-1=0\); thus, since the quantities \(1\ (=r^0),\allowbreak r,\allowbreak rr\ldots r^{n-1}\) are all distinct, they give all the roots of the equation \(x^n-1=0\), and hence the collection \(r,\allowbreak rr,\allowbreak r^3\ldots r^{n-1}\) coincides with \(\Omega\). More generally, it follows easily that \(\Omega\) coincides with the collection \(r^e,\allowbreak r^{2e},\allowbreak r^{3e}\ldots r^{(n-1)e}\) if \(e\) is any integer not divisible by \(n\), positive or negative. We therefore have \[X=(x-r^e)(x-r^{2e})(x-r^{3e})\ldots(x-r^{(n-1)e})\] whence \[\begin{gathered}
r^e+r^{2e}+r^{3e}+\ldots+r^{(n-1)e}=-1,\\
\text{and}\quad1+r^e+r^{2e}+\ldots+r^{(n-1)e}=0
\end{gathered}\]

We will call two roots such as \(r\) and \(\frac1r\ (=r^{n-1})\), or generally \(r^e\) and \(r^{-e}\), \textit{reciprocal}; clearly the product of the two simple factors \(x-r\) and \(x-\frac1r\) is real, \(=xx-2x\cos\omega+1\), where the angle \(\omega\) is equal either to the angle \(\frac Pn\) or to one of its multiples.

\clearpage
\Needspace{6\baselineskip}\Article{340}
Thus, if one root in \(\Omega\) is denoted by \(r\), all the roots of the equation \(x^n-1=0\) are expressed by powers of \(r\), and a product formed in any way from several roots of this equation can be represented by \(r^\lambda\), where \(\lambda\) is either \(0\), or positive and \(<n\). Now let \(\varphi(t,\allowbreak u,\allowbreak v\ldots)\) denote a rational integral algebraic function of the indeterminates \(t,\allowbreak u,\allowbreak v\), etc., expressible as a sum of terms of the form \(ht^\alpha u^\beta v^\gamma\ldots\). It is clear that, if we substitute for \(t,\allowbreak u,\allowbreak v\), etc., some of the roots of the equation \(x^n-1=0\), say \(t=a,\allowbreak u=b,\allowbreak v=c\), etc., \(\varphi(a,\allowbreak b,\allowbreak c\ldots)\) can be put in the form \[A+A^{\prime} r+A^{\prime\prime} rr+A^{\prime\prime\prime} r^3+\ldots+A^\nu r^{n-1}\] where the coefficients \(A,\allowbreak A^{\prime}\), etc. (some of which may be absent and hence \(=0\)) are fixed quantities; moreover, all these coefficients will be integers if all the fixed coefficients in \(\varphi(t,\allowbreak u,\allowbreak v\ldots)\), that is, all the
\SourceLeaf{422}
\(h\), are integers. If we then replace \(t,\allowbreak u,\allowbreak v\ldots\) by \(aa,\allowbreak bb,\allowbreak cc\ldots\), respectively, any term such as \(ht^\alpha u^\beta v^\gamma\ldots\), which previously reduced to \(r^\sigma\), will now become \(r^{2\sigma}\). It follows easily that \[\varphi(aa,bb,cc\ldots)=A+A^{\prime} rr+A^{\prime\prime} r^4+A^{\prime\prime\prime} r^6+\ldots+A^\nu r^{2n-2}\] Similarly, for any integer value of \(\lambda\), we have in general \[\varphi(a^\lambda,b^\lambda,c^\lambda\ldots)=A+A^{\prime} r^\lambda+A^{\prime\prime} r^{2\lambda}+\ldots+A^\nu r^{(n-1)\lambda}\] This proposition is of the greatest importance and forms the foundation of the following investigations. — This also gives \[\varphi(1,1,1\ldots)=\varphi(a^n,b^n,c^n\ldots)=A+A^{\prime}+A^{\prime\prime}+\ldots+A^\nu\] and \[\begin{aligned}
\varphi(a,b,c\ldots)+\varphi(aa,bb,cc\ldots)&+\varphi(a^3,b^3,c^3\ldots)\\
&+\ldots+\varphi(a^n,b^n,c^n\ldots)=nA
\end{aligned}\] Thus this sum is always an integer divisible by \(n\) whenever all the fixed coefficients in \(\varphi(t,\allowbreak u,\allowbreak v\ldots)\) are integers.

